\begin{claim}{D.2.46}{Proved}{The universal Moore composite lies in one cyclic subgroup}
Retain the two-local notation of D.2.44, and put
\[
 h=h_Ml_4\in\pi_{14}Y,\qquad
 e=u_4E\epsilon',\qquad
 q=\alpha\sigma'+e,\qquad \alpha=u_4\nu_4.
 \tag{D.2.46a}
\]
Then
\[
             4q=0,\qquad h\in\langle q\rangle,
             \qquad 2h=2q.
 \tag{D.2.46b}
\]
If $q$ has order four, then $h=q$ or $h=-q$.  If $q$ has
order at most two, then $h=0$ or $h=q$.  More precisely, with
compatible sphere-generator signs, the universal composite itself
satisfies
\[
 wl_4\in
 \left\langle i_{5*}\bigl(\nu_4\sigma'
                    -\varepsilon E\epsilon'\bigr)\right\rangle,
 \qquad \varepsilon\in\{1,-1\}.
 \tag{D.2.46c}
\]
This sharpens the two-component span in D.2.44.  It does not
determine the order of $q$ in the full target or choose between
the possibilities in (D.2.46b).
\end{claim}
\begin{proof}
The stronger suspension constraint is obtained in the actual
two-cell complex $A_5$, before mapping to $P_5$.  This order of
the argument is essential: D.2.44 proves that suspension in
$P_5$ alone cannot distinguish the two candidate components.

Use the actual cofiber map $c:P_8\to A_5$ and the characteristic
class $\tau_5$ from (D.2.44g)--(D.2.44n).  In particular,
\[
 A_5=S^4\cup_{2\omega_4}e^8,
 \qquad \omega_4=[\iota_4,\iota_4],
 \qquad c|_{S^7}=\omega_4.
 \tag{D.2.46d}
\]
Write $i:S^4\to A_5$ for its bottom inclusion and $j_*$ for
the map to the relative group of $(A_5,S^4)$.  The mixed relative
Jacobi calculation of (D.2.44n), in this pair itself, gives
\[
 j_*\bigl(3cl_4-2[\tau_5,i]\bigr)=0.
\]
Exactness therefore supplies a sphere class
$\delta_0\in\pi_{14}S^4_{(2)}$ with
\[
             i_*\delta_0=3cl_4-2[\tau_5,i]
                       \quad\text{in }\pi_{14}A_5.
 \tag{D.2.46e}
\]
There is no inference here from an equality only after the
twelve-cell has been attached.

Both terms on the right of (D.2.46e) suspend to zero.  This is
the usual suspension property for the Whitehead product.  For
the other term, the bottom restriction of
$\Sigma c:P_9\to\Sigma A_5$ is $E\omega_4=0$.
The Moore Puppe sequence makes $\Sigma c$ factor through the
actual pinch $p_9:P_9\to S^9$.  The Stiefel attachment has
\[
                 El_4=i_9[\iota_8,\iota_8],
 \tag{D.2.46f}
\]
as in \cite[Lemmas 2.7--2.8, p.~523]{EH5MorisugiMukai2000}.
Its composite with $p_9$ is zero, so $E(cl_4)=0$.
Finally the suspended attaching map of $A_5$ is zero, giving
\[
          \Sigma A_5\simeq S^5\vee S^9
          \quad\text{under the bottom }S^5.
\]
The bottom inclusion thus has a retraction.  Suspending
(D.2.46e) and applying this retraction proves the stronger
sphere-level equality
\[
                 E\delta_0=0\quad\text{in }\pi_{15}S^5_{(2)}.
 \tag{D.2.46g}
\]

For completeness, the entire sphere suspension kernel is explicit.
The primary tables and suspension formulas reproduced in
\cite[Proposition 7 and its proof, p.~7]{FGHStricklandToda}
give
\[
\begin{aligned}
 \pi_{14}S^4_{(2)}
  &=\mathbb Z/8\{a=\nu_4\sigma'\}
    \oplus\mathbb Z/4\{b=E\epsilon'\}
    \oplus\mathbb Z/2\{c_0=\eta_4\mu_5\},\\
 \pi_{15}S^5_{(2)}
  &=\mathbb Z/8\{d=\nu_5\sigma_8\}
    \oplus\mathbb Z/2\{e_0=\eta_5\mu_6\},\\
 Ea&=2d,\qquad Eb=\varepsilon\,2d,\qquad Ec_0=e_0,
 \qquad \varepsilon\in\{1,-1\}.
\end{aligned}
 \tag{D.2.46h}
\]
These are the stated Toda Theorem 7.3 and equations (7.10),
(7.16) inputs in that proof; no general Hopf-composition formula
is needed for this kernel calculation.  A class $ma+nb+tc_0$
belongs to the kernel precisely when $t=0$ and
$m+\varepsilon n\equiv0\pmod4$.  Hence
\[
 \ker\bigl(E:\pi_{14}S^4_{(2)}\to\pi_{15}S^5_{(2)}\bigr)
       =\mathbb Z/8\{a-\varepsilon b\}.
 \tag{D.2.46i}
\]

The actual third cell of the pinch fiber kills $2\tau_5$.
Moreover the full Hopf--Hilton comparison in D.2.44 proves
$wl_4=(A_5\to P_5)cl_4$, with its mixed correction retained
and shown to vanish.  Mapping (D.2.46e) to $P_5$ therefore gives
\[
                        3wl_4=i_{5*}\delta_0.
 \tag{D.2.46j}
\]
Equations (D.2.46g)--(D.2.46i) and invertibility of three on a
two-primary cyclic group prove (D.2.46c).

In $Y$, D.2.44 gives $2e=0$, so the sign of the $e$ term
does not affect the generated subgroup.  It remains to justify
the exponent and the double without discarding a Hilton term.
The sphere $S^7$ is an H-space, and consequently
\[
 [\alpha,\alpha]=\alpha_*[\iota_7,\iota_7]=0.
\]
For precomposition by $\sigma':S^{14}\to S^7$, the only
possible nonlinear Hilton term is a multiple of this
Whitehead square: the two-letter factor begins in dimension
thirteen, while every three-letter factor begins in dimension
nineteen.  Thus precomposition by $\sigma'$ is additive on
the cyclic subgroup generated by $\alpha$.  The already
proved order-four relation for $\alpha$ in D.2.18 yields
\[
              4(\alpha\sigma')=(4\alpha)\sigma'=0,
              \qquad 4q=0.
 \tag{D.2.46k}
\]
Finally the exact Stiefel double from D.2.41 is
\[
              2h=(-2\alpha)\sigma'
                 =-2(\alpha\sigma')=2q.
 \tag{D.2.46l}
\]
The last equality uses (D.2.46k) and $2e=0$.
Together with $h\in\langle q\rangle$, this proves
(D.2.46b) and the two stated order-dependent alternatives.
In particular no unconditional odd-coefficient assertion is
being made when the double vanishes.
\end{proof}
